\documentclass[a4paper,10pt]{article}

% --- PACKAGES ---
\usepackage[utf8]{inputenc}
\usepackage[T1]{fontenc}
\usepackage[scaled]{helvet}
\renewcommand{\familydefault}{\sfdefault}
\usepackage[left=1.4cm,right=1.4cm,top=1.2cm,bottom=1.2cm]{geometry}
\usepackage{enumitem}
\setlist{leftmargin=*,noitemsep,topsep=0.05em,partopsep=0em,parsep=0em}
\usepackage{parskip}
\setlength{\parskip}{0.1em}
\setlength{\parindent}{0pt}
\setlength{\emergencystretch}{3em}
\hfuzz=2pt
\usepackage{xcolor}
\definecolor{myblue}{HTML}{205792}
\usepackage[colorlinks=true,urlcolor=myblue,linkcolor=myblue,citecolor=myblue]{hyperref}
\usepackage{microtype}
\newcommand{\cventry}[2]{\noindent\begin{minipage}[t]{\linewidth}\textbf{#1}\hfill\textbf{#2}\end{minipage}}
\hypersetup{pdftitle={Renato Mignone CV},pdfauthor={Renato Mignone}}
% \renewcommand{\href}[2]{\textcolor{myblue}{#2}}   % <-- add this
\pagestyle{empty}
\setcounter{secnumdepth}{0}


% --- SECTION TITLE STYLE ---
% \newcommand{\mysection}[1]{%
%   \vspace{0.3ex}
%   {\large\bfseries\scshape\raggedright\color{myblue} #1}
%   \rule{\linewidth}{0.4pt}
%   \vspace{-0.3em}
% }

\newcommand{\mysection}[1]{%
  \vspace{0.3ex}
  {\large\bfseries\scshape\color{myblue} #1}\enspace%
  \leaders\hrule height 0.4pt\hfill\kern0pt
  \vspace{-0.2em}
}



\begin{document}

% --- HEADER ---
\begin{center}
    {\LARGE\scshape\color{myblue} Renato Mignone}\\[0.15em]
    {\normalsize Cybersecurity Engineer | Post-Quantum Cryptography, AI Security}\\[0.25em]
    \href{tel:+391234567890}{+39 123 456 7890} \quad
    \href{mailto:renato.mignone@gmail.com}{renato.mignone@gmail.com} \quad
    \href{https://www.linkedin.com/in/renato-mignone}{linkedin.com/in/renato-mignone} \quad
    \href{https://github.com/RenatoMignone}{github.com/RenatoMignone}\\
    \href{https://renatomignone.github.io}{renatomignone.github.io} \enspace|\enspace Düsseldorf, Germany\enspace|\enspace EU Work Authorization\enspace|\enspace Open to Relocation\enspace|\enspace Available Sept 2026
\end{center}

% --- TECHNICAL SKILLS ---
\mysection{Technical Skills}
\begin{itemize}
    \item \textbf{Security Engineering:} Penetration Testing, Vulnerability Research \& Exploitation, Kernel Security \& eBPF, Binary Analysis (Ghidra, GDB, radare2), Malware Analysis, Web Application Security (OWASP Top 10)
    \item \textbf{Threat Detection \& Response:} MITRE ATT\&CK Framework, Threat Modeling, Network Intrusion Detection (NIDS/IDS), Anomaly Detection, Incident Response, Threat Intelligence, Detection Engineering
    \item \textbf{Cryptography:} Post-Quantum Cryptography (Lattice-Based, Hash-Based), Blind Signatures, Zero-Knowledge Proofs (ZKP), OPRFs, Anonymous Tokens, PKI, OpenSSL, Cryptographic Protocol Design
    \item \textbf{AI for Cybersecurity:} AI-Based Threat Detection, PyTorch, Scikit-learn, GNN, RNN, LSTM, FFNN, Transformers (BERT, UniXcoder), Anomaly Detection Models, Adversarial ML
    \item \textbf{Programming:} Python (multithreading, automation, ML/AI), C/C++ (systems, cryptographic implementations), Bash, Assembly (RISC-V, ARM), SQL, Java, JavaScript/TypeScript, Rust, Node.js/Express, React, PostgreSQL
    \item \textbf{Systems \& Platforms:} Linux (Kali, Ubuntu), Docker, Container Security, QEMU, Git, bpftool, Wireshark, Nmap, GNS3
\end{itemize}

% --- EXPERIENCE ---
\mysection{Experience}
\cventry{Security Engineer Intern, Huawei - Data Privacy Lab, D\"usseldorf, Germany}{Feb 2026 - Present}
\begin{itemize}
    \item Engineered a PBAT (Policy-Based Anonymous Token) layer on a STARK-based post-quantum circuit, extending the trace by 105 Poseidon columns (+15\%) and 210 degree-3 transition constraints while preserving 115-bit security.
    \item Benchmarked the full construction end-to-end: proof generation 294\,ms to 520\,ms, verification 29.2\,ms to 32.3\,ms (+10.6\%), proof size 172.7\,KB to 191.5\,KB; all 4 integration tests passing.
    \item Analyzed cryptographic protocols for privacy-preserving tokens, including PQ blind signatures, zero-knowledge proofs, and OPRFs; evaluated CRYSTALS-Dilithium, SPHINCS+, and PS signatures under formal game-based security models.
\end{itemize}
% \begin{itemize}
%     \item Researching post-quantum secure policy-based anonymous authentication tokens, exploring lattice-based and hash-based cryptographic constructions resistant to quantum attacks.
%     \item Analyzing cryptographic protocols including blind signatures, zero-knowledge proofs, and oblivious pseudorandom functions for privacy-preserving token systems.
%     \item Implementing and evaluating constructions using CRYSTALS-Dilithium, Sphincs+, and PS signatures with focus on formal security models (game-based and simulation-based proofs).
% \end{itemize}

\mysection{Leadership}
\cventry{Logistics and Fundraising, IEEE-HKN (Mu Nu Chapter) - Honor Society}{Apr 2024 - Present}
\begin{itemize}
    \item Coordinated 5+ STEM events and hackathons (150-250 participants): logistics, infrastructure, and support.
    \item Secured sponsorships and hardware from ARM, NXP, and STMicroelectronics for university events.
    \item Built secure full-stack app (React/Node.js/PostgreSQL, TOTP 2FA, JWT, RBAC) - 1st place, IEEE-HKN Hackathon 2025.
\end{itemize}

% --- PROJECTS ---
\mysection{Key Projects}
\vspace{0.2em}

\cventry{\textbf{AI-Driven Threat Detection Suite} - \href{https://github.com/RenatoMignone/AI-Driven-Threat-Detection-Research}{GitHub}}{Jan 2026}
\begin{itemize}
    \item Multi-model defense system leveraging FFNNs, RNNs, GNNs to detect malware, network intrusions, and bash-level attacks.
    \item Achieved 98.1\% accuracy in malware classification by modeling API call sequences with GNNs and RNNs.
    \item Developed an unsupervised network IDS with 95\% precision using deep autoencoders for anomaly detection.
    \item Built an NLP-based MITRE ATT\&CK tactic classifier to map bash command sessions using BERT and UniXcoder.
    \item \textbf{Stack:} PyTorch, GNN, RNN, LSTM, FFNN, Transformers (BERT, UniXcoder), Deep Autoencoders, Malware Analysis
\end{itemize}

\vspace{0.2em}
\cventry{\textbf{Linux Kernel eBPF Verifier Security Analysis} - \href{https://github.com/RenatoMignone/Linux-eBPF-Verifier-Bypass-Research}{GitHub}}{Sept 2025}
\begin{itemize}
    \item Conducted security analysis of Linux kernel 6.8 eBPF verifier, identifying privilege escalation vulnerabilities based on ISO-IEC TS 17961-2013 secure coding standards.
    \item Developed proof-of-concept exploits demonstrating kernel privilege escalation through verifier logic bypass techniques (5 confirmed exploitable vulnerabilities from 60+ candidates).
    \item Performed multi-stage vulnerability analysis covering C source code, eBPF bytecode verification, JIT compilation semantics, and runtime exploitation using QEMU, bpftool, and custom static analysis tools.
    \item \textbf{Stack:} Kernel Exploitation, eBPF Internals, Vulnerability Research, Reverse Engineering, C/Assembly, QEMU, GDB
\end{itemize}

\vspace{0.2em}
\cventry{\textbf{Goodware Equivocal Behavior Analysis System} - (private - research project)}{Sept 2024}
\begin{itemize}
    \item Designed and implemented an asynchronous, multithreaded Python CLI tool integrating 3 sandbox platforms (Hybrid Analysis, VirusTotal, ANY.RUN) to analyze 36 trusted binaries in parallel across Windows 10 and Windows 11 environments.
    \item Defined 12 Equivocal Software Behaviours from 60 MITRE ATT\&CK Enterprise Matrix techniques; detected System Reconnaissance and Advanced OS Utility Exploitation in 100\% of analyzed systems, including major browsers.
    % \item Produced per-binary, per-sandbox structured JSON reports and post-processed results with Pandas and Seaborn to quantify undisclosed covert reconnaissance activity in goodware, contributing to the software transparency research agenda.
    \item \textbf{Stack:} Python, Multithreading, REST APIs, MITRE ATT\&CK, Hybrid Analysis, VirusTotal, ANY.RUN, Supply Chain Security
\end{itemize}


% --- EDUCATION ---
\mysection{Education}
\cventry{MSc in Cybersecurity Engineering\enspace|\enspace Politecnico di Torino, Italy}{Sept 2024 - Oct 2026}\\
GPA: 29.48/30 (top 5\% of cohort) \enspace|\enspace English-taught, ECSO-aligned curriculum\\
Thesis: \textit{Post-Quantum Secure Policy-Based Anonymous Authentication Tokens} (Huawei Data Privacy Lab)\\
% Relevant Coursework: Advanced Cryptography, Security Verification and Testing, Network and Cloud Security, Hardware and Wireless Security, AI for Cybersecurity, Advanced Information Systems Security
\vspace{0.15em}
\cventry{BSc in Computer Engineering\enspace|\enspace University of Sannio, Italy}{Sept 2021 - June 2024}\\
Grade: 101/110 \enspace|\enspace Italian-taught curriculum\\
Thesis: \href{https://zenodo.org/records/19262869}{\textit{Extracting Equivocal Behaviors from Trusted Systems}} (system integrity \& anomaly detection)

% --- CERTIFICATIONS, LEADERSHIP & LANGUAGES ---
\mysection{Certifications, Achievements \& Languages}

\begin{itemize}
    \item \textbf{Achievements:} Winner (1st place), IEEE-HKN International Hackathon 2025; Top 60 (57th/2000 Teams), Reply AI Agent Challenge 2026; Top 100 (88th), Reply Hack The Code Challenge 2025 (56,196,106 points).
    \item \textbf{Certifications:} GNU/Linux Advanced (NetStudent \& WEEE Open, 2025); Cambridge English B2 First (FCE, 2024).
    \item \textbf{Languages:} Italian (Native), English (Professional Working Proficiency), German (Basic).
\end{itemize}

\end{document}